\centerline{\bf \Huge Планиметрия}



\problem{\textit{\textbf{1}}} Основания равнобедренной трапеции равны 8 и 18, а периметр равен 56.

Найдите площадь трапеции.

\problem{\textit{\textbf{2}}} Окружности с центрами в точках P и Q не имеют общих точек, и ни одна из них не лежит внутри другой. Внутренняя общая касательная к этим окружностям делит отрезок, соединяющий их центры, в отношении a:b. 

Докажите, что диаметры этих окружностей относятся как a:b.

\problem{\textit{\textbf{3}}} В трапеции ABCD боковая сторона AB перпендикулярна основанию BC. Окружность проходит через точки C и D и касается прямой AB в точке E. 

Найдите расстояние от точки E до прямой CD, если AD = 14, BC = 12.


\problem{\textit{\textbf{4}}} В треугольнике ABC на его медиане BM отмечена точка K так, что BK : KM = 7 : 3. Прямая AK пересекает сторону BC в точке P.

Найдите отношение площади треугольника BKP к площади четырехугольника KPCM.

\problem{\textit{\textbf{5}}} Четырехугольник ABCD с перпендикулярными диагоналями AC и BD вписан в окружность.

а) Докажите, что прямая, проходящая через точку пересечения диагоналей четырехугольника перпендикулярно стороне BC, делит пополам сторону AD.

б) Найдите стороны четырехугольника ABCD, если известно, что AC=84 и BD=77, а диаметр окружности равен 85.

\problem{\textit{\textbf{6}}} В трапеции ABCD точка E — середина основания AD, точка M — середина боковой стороны AB. Отрезки CE и DM пересекаются в точке O.

а) Докажите, что площади четырёхугольника AMOE и треугольника COD равны.

б) Найдите, какую часть от площади трапеции составляет площадь четырёхугольника AMOE, если BC=3, AD=4.

